%==============================================================================
% Appendix: Datasheet for Datasets (Gebru et al., 2021)
% Standalone section file. Included after the main body inside the \appendix
% block opened by main.tex (do NOT re-issue \appendix here, or the appendix
% counter resets and the datasheet is renumbered ahead of the other appendices).
%==============================================================================

\section{Datasheet for Datasets}
\label{app:datasheet}

Following the documentation framework of Gebru et al.~\citep{gebru2021datasheets},
we answer the standard datasheet questions for the \bench{} corpus. The seven
labeled subsections below cover motivation, composition, the collection and
generation process, preprocessing and cleaning (including a complete known-errors
list), intended uses, distribution, and maintenance. All counts are taken
directly from the released artifacts on disk.

%------------------------------------------------------------------------------
\subsection{Motivation}
%------------------------------------------------------------------------------

\textbf{For what purpose was the dataset created?}
\bench{} was created to measure \emph{long-horizon organizational memory} in AI
agents: the ability to answer everyday questions about an organization
(``when did we change enterprise pricing?'', ``who owns the Acme relationship?'',
``what did engineering decide about the migration last quarter, and has that
changed?'') by stitching evidence across channels, across people, and across
time. Existing memory benchmarks (LoCoMo, LongMemEval, MemBench,
MemoryAgentBench, HaluMem) are conversational or personal-scale, cap at roughly
$200\text{K}$--$1\text{M}$ tokens of mostly two-party dialogue, and are already
saturating, with leading systems exceeding $90\%$. No public benchmark measures
the organizational regime: multi-source, multi-author, multi-channel corpora at
the scale a company accumulates over years, where the answer to a question may
have been superseded, corrected, or never written down explicitly. \bench{}
targets precisely that gap.

\textbf{Who created the dataset and on whose behalf?}
The dataset was created by the authors of this paper as part of the \bench{}
benchmark effort.

\textbf{Who funded the creation of the dataset?}
Generation used a single locally served open-weight model
(\texttt{gpt-oss:20b}, served via Ollama at an internal endpoint), so corpus
generation incurred no external API spend; evaluation of the four public memory
systems used the authors' own vendor accounts and API credits. No external grant
funded the dataset.

%------------------------------------------------------------------------------
\subsection{Composition}
%------------------------------------------------------------------------------

\textbf{What do the instances represent?}
The dataset describes a single, fully fictional organization, \emph{Helix
Logistics}, a small-to-mid SaaS/freight-tech company with a multi-year narrative
arc (founded 2020; Helix API in 2021; Series A and a B2C$\rightarrow$mid-market
pivot in 2022; a 2023 customer escalation, layoff, and COO change; Helix AI
Optimizer in 2024; a documentation overhaul in 2025; a planned Series B in 2026).
There are two instance types. (i)~\emph{Corpus artifacts}: individual
documents/threads from one channel (Slack threads, email threads, meeting notes,
meeting transcripts, incident reports, postmortems, ADRs, Notion docs, and CRM
entries), each carrying a structured header (\texttt{slot\_id},
\texttt{event\_id}, genre, author persona, date, witness \texttt{role}, and
\texttt{testimony\_type}). (ii)~\emph{Questions}: natural-language queries with a
category-specific structured ground-truth answer, weighted rubric subpoints,
evidence-artifact citations, and metadata. A separate \texttt{personas.jsonl}
registry and a \texttt{source\_of\_truth/\{events,relations\}.jsonl} graph
support both.

\textbf{How many instances are there in total?}
\bench{} ships three scale tiers of the Helix corpus. The \textbf{small} tier
contains $17$ events ($121$ corpus artifacts) and $11$ questions. The
\textbf{medium} tier contains $157$ events and $443$ corpus artifacts and $73$
questions. The \textbf{large} tier contains $113$ questions over a
multi-million-token corpus. Per-genre artifact counts and the role
(primary/secondary witness) and testimony (\texttt{direct}/\texttt{inferred})
splits are tabulated in Appendix~\ref{sec:design:personas}; most facts must be
triangulated from secondary, partly inferential sources rather than read off a
single canonical record. The persona registry holds $72$ personas, each with a
dated role history, voice markers, signature phrases, channel-mix probabilities,
and founder / long-tail flags. XL and XXL tiers, and additional companies of
other sizes and verticals, are planned but not part of this release.

\textbf{Does each instance contain labels or ground truth?}
Yes. Every question carries a category-specific structured
\texttt{ground\_truth\_answer} that is a deterministic projection of the
source-of-truth graph (not a string recalled by the generating model), plus
weighted \texttt{rubric\_subpoints}; the full per-category field schema
(C1--C6) is given in Appendix~\ref{sec:design:rubric}. The small and medium
tiers cover exactly these six
categories (C1--C6). A further EMERGENT category (de-facto patterns never stated
as an explicit rule, with multi-facet answers) is evaluated only at the large
tier and above, where a deep behavioural corpus is needed for the patterns to be
reconstructable; it is out of scope for the small and medium tiers in this paper.
Each question also carries three paraphrases of its text, and the
\texttt{rubric\_subpoints} above are weighted for facet-coverage scoring.

\textbf{Is any information missing from individual instances?}
A small number of citations are known to be unresolvable and are flagged in the
errata (see \emph{Preprocessing/Cleaning} below); these are documented rather
than silently dropped. The EMERGENT question category is out of scope for the
small and medium tiers in this paper and is evaluated only at the large tier and
above; the small corpus still physically retains its emergent-evidence background
documents (kept as benign background noise), but the emergent \emph{questions}
were removed from the small tier, so no EMERGENT questions appear in either the
small or medium release here.

\textbf{Are relationships between instances made explicit?}
Yes. Corpus artifacts are grouped into event folders (\texttt{EV-YYYY-NNN/}),
multiple artifacts may witness the same event from different angles and channels,
and each question's \texttt{evidence\_artefact\_ids} field ties its answer to the
specific artifacts that license it. The \texttt{source\_of\_truth} graph encodes
bi-temporal relations (supersessions, corrections, contradictions) between facts.

\textbf{Are there recommended data splits?}
The dataset is structured as difficulty/scale tiers (small / medium / large)
rather than train/validation/test splits, because it is an evaluation-only
benchmark: there is no training set. We report results per tier.

\textbf{Is the dataset self-contained?}
Yes. The full Helix corpus, ground truth, persona registry, and source-of-truth
graph are vendored in-repo under \texttt{datasets/helix/\{small,medium,large\}/}.
It does not link to or depend on any external resource.

\textbf{Does the dataset contain confidential or offensive content, or data
relating to identifiable people?}
No. The organization, its employees, customers, events, and all artifact text
are entirely fictional and machine-generated. No real company's communications
and no real individuals' data appear in the dataset, so there is no personally
identifiable information and no confidentiality risk.

%------------------------------------------------------------------------------
\subsection{Collection / Generation Process}
%------------------------------------------------------------------------------

\textbf{How was the data acquired?}
The Helix corpus is \emph{entirely synthetic}; no human wrote any Helix email,
Slack message, transcript, ADR, or ticket. Every artifact was produced by
parameterized prompts sent to a single locally served open-weight model
(\texttt{gpt-oss:20b}). A six-stage pipeline (\texttt{helix-corpus}) runs in a
fixed order so that ground truth is fixed \emph{before} any natural-language text
is generated:

\begin{itemize}
  \item \textbf{Stage A --- Company canon.} Instantiates the fictional company:
    founders, multi-year growth arcs, an $80$--$90$ person persona registry with
    dated roles, customer arcs, documentation-discipline eras, and a tooling
    timeline with explicit replacements (e.g., Salesforce replaces HubSpot in
    2022-Q1; Confluence replaces Notion in 2023-Q1). This is the only layer a
    human configured, by choosing scale parameters (company size, year range,
    team topology).
  \item \textbf{Stage B --- Source-of-truth graph.} Generates the ground-truth
    event log and bi-temporal relation graph (supersessions, corrections,
    contradictions). All downstream answers are projected deterministically from
    this graph; the model never invents a ground truth during question
    generation.
  \item \textbf{Stage C --- Corpus.} Generates a cluster of two to three
    artifacts per event across the genre taxonomy, each with a structured header.
  \item \textbf{Stage D --- Cross-references and noise injection.} Adds
    cross-artifact references and realistic distractor content (drafts,
    half-finished threads, contradictions among secondary witnesses), recording
    every injected distractor in \texttt{noise\_manifest.jsonl}.
  \item \textbf{Stage E --- Question generation.} Produces questions whose
    answers are derived from the Stage B graph; the model's role is phrasing,
    not answer invention.
  \item \textbf{Stage F --- Self-consistency validation.} Reviews output for
    answerability, rubric coherence, and distractor integrity, emitting a
    \texttt{validation\_report.jsonl}; passing questions become
    \texttt{benchmark\_v0.0.jsonl}.
\end{itemize}

\textbf{What is the bi-temporal model?}
Every event/fact in the source-of-truth graph carries four timestamps:
\texttt{valid\_at} (became true in the world), \texttt{invalid\_at} (stopped
being true), \texttt{ingested\_at} (recorded by the organization), and
\texttt{expired\_at} (record retracted or corrected). Facts are invalidated,
never deleted, preserving a no-loss audit trail; supersession links connect an
old fact to its replacement with an explicit reason; and questions can ask what
the organization believed \texttt{as\_of} a past date, distinct from what is
true now. The harness is capability-aware: a system that by design cannot answer
as-of-date temporal queries is reported N/A for that category rather than
silently scored zero.

\textbf{Were any ethical review processes conducted?}
No human-subjects review was required: the data is fully synthetic and contains
no real individuals or organizations.

\textbf{Known methodological limitation (circularity).}
Single-model end-to-end generation is circular by construction: the same model
writes the artifacts, the questions, and the answers, and validates them.
Structural defenses (closed-enum event types, graph-projected rather than
model-recalled answers, multi-pass critique, and automated sentinels for persona
consistency and anachronism) reduce but do not eliminate this circularity. The
README acknowledges this explicitly, and a cross-model validation pass (a
different model family as validator) is planned for a future revision.

%------------------------------------------------------------------------------
\subsection{Preprocessing / Cleaning / Labeling}
%------------------------------------------------------------------------------

Because ground truth is projected from the source-of-truth graph rather than read
back out of free text, label ambiguity is bounded by construction. Post-generation
audits nonetheless identified six classes of defects, corrected across three audit
waves. All corrections are logged in the per-tier \texttt{CHANGELOG.md} files, and
pre-patch snapshots are preserved (\texttt{benchmark\_v0.0.original.jsonl} for the
small tier; \texttt{benchmark\_v0.0.pre-wave2.jsonl} and
\texttt{benchmark\_v0.0.pre-wave3.jsonl} for the medium tier), so the audit trail
is complete. We list the known errors and corrections in full.

\begin{enumerate}
  \item \textbf{Slot-ID collisions (medium tier).} The Stage C corpus generator
    reset its per-batch slot-ID counter to \texttt{-001} while the event-ID
    counter kept advancing, producing $443$ corpus files with only $311$ unique
    slot IDs ($24$ slot IDs appeared $2$--$10$ times across different event
    directories). A wave-1 audit renamed $140$ files on disk and in
    \texttt{corpus\_index.jsonl} (rebuilding the index from scratch) and
    repointed $35$ question rows whose \texttt{evidence\_artefact\_ids} had
    resolved to the wrong event cluster. Three slot IDs
    (\texttt{ART-EV-2023-002-003}, \texttt{ART-EV-2023-006-003},
    \texttt{ART-EV-2023-066-001}) remained unresolvable because the
    home-directory artifacts were never generated; those citations were left
    intact and flagged.
  \item \textbf{Persona registry gaps and surname cross-contamination.} The
    generator occasionally fused names across personas (e.g., ``Daniel Patel''
    for Daniel Ortiz, ``Arjun Patel'' for Arjun Mehta, ``Ethan Patel'' for Ethan
    Brooks, ``Miguel Alvarez'' for Miguel Torres, ``Luis Hernandez'' for Luis
    Ramirez). In the small tier, two artifacts referred to a ``Maya Singh'', a
    generation-time typo for Maya Patel (CEO), confirmed by the absence of any
    ``Maya Singh'' in the registry and the event date predating the only other
    Maya. The medium tier had two duplicate display-name collisions (two distinct
    ``Daniel Kim'' personas; two distinct ``Marco Rossi'' personas); the
    lower-seniority duplicate in each pair was renamed (\texttt{Daniel R. Kim},
    \texttt{Marco S. Rossi}). Four personas referenced in artifacts but absent
    from the registry (P-0009 Omar Al-Jabri, P-0090 Sarah Patel, P-0091 Ethan
    Kim, P-0092 Emily Carter) were added as stubs. Four ambiguous name
    occurrences (Marco Ruiz, Kevin Patel, Maya Lin, Luis Martinez) lacked
    sufficient context for confident assignment and were left unchanged with a
    flag for future resolution.
  \item \textbf{Ground-truth-to-corpus label drift (testimony types).}
    {\sloppy
    Generator-internal labels (\texttt{"direct\_testimony"},
    \texttt{"inference"}) diverged from the corpus header convention
    (\texttt{"direct"}, \texttt{"inferred"}). Ninety field values in the medium
    tier were normalized ($45$ in \texttt{testimony\_attribution} and $45$ in
    \texttt{evidence\_summary} strings across $15$ C5 rows); the same fix had
    been applied earlier to the small tier.\par}
  \item \textbf{Phantom artifact citations.} Several questions in both tiers
    cited \texttt{evidence\_artefact\_ids} pointing to slot IDs with no file on
    disk. In the small tier, two primary artifacts (a postmortem for EV-2022-008
    and a meeting transcript for EV-2022-013) were entirely absent; both were
    written and added. In the medium tier, two questions' citation lists were
    repointed post wave-3 to the correct anchor events after the slot-ID
    collision fix had made the original citations unresolvable.
  \item \textbf{Spurious supersession entries and C4 audit-replay integrity
    (small tier).} A wave-1 audit found that some C4 audit-replay questions cited
    artifacts that did not exist, and that one question's supersession structure
    was over-specified into a negative-knowledge test with an empty
    \texttt{replacements=[]}. Phantom-citation removals were applied to five
    questions; one question (Q-0010) was rewritten to reward the correct
    bi-temporal supersession pattern, and another (Q-0011) had its
    \texttt{resolution\_status} corrected and a rubric criterion inverted.
  \item \textbf{Natural-language question rewording (small tier).} Eleven small-tier
    question texts (of the question set as it then stood, prior to deferring the
    emergent-pattern category to the large tier) contained benchmark-design
    artifacts that
    would not appear in natural user queries (``the system's understanding'',
    ``your statement'', ``your claim'', ``reconstruct the knowledge held by the
    system''). These were reworded to phrasings a person would actually use. All
    ground-truth answers and rubric sub-criteria were left unchanged; the prior
    text for each modified question is preserved at
    \texttt{metadata.text\_pre\_rewording\_2026\_05\_26}.
\end{enumerate}

\textbf{Is the raw (pre-cleaning) data available?}
Yes. Pre-patch snapshots are preserved at every audit wave (filenames above), so
the raw generated output and each intermediate state remain inspectable alongside
the corrected release.

%------------------------------------------------------------------------------
\subsection{Uses}
%------------------------------------------------------------------------------

\textbf{What tasks is the dataset intended for?}
\bench{} is an evaluation-only benchmark for long-horizon organizational memory
in AI agents and memory systems. Its intended use is to measure how well a memory
system can ingest a multi-source, multi-author, multi-channel organizational
corpus and then answer questions requiring supersession reasoning, decision
provenance, bi-temporal (as-of-date) retrieval, audit replay, justification
chains, and contradiction resolution. Because the tiers form a scaling ladder,
the dataset is specifically intended to expose the \emph{scaling curve}: how a
system degrades as the corpus grows from conversational scale to company scale.

\textbf{Has the dataset been used for any tasks already?}
Yes. In this paper we evaluate four publicly available memory systems on the
small tier and report a medium-tier result for the strongest of them, under a
uniform, capability-aware harness with rubric-based scoring.

\textbf{Is there anything that should not be done with the dataset?}
The dataset must not be used as \emph{training} data for any system that will be
evaluated on it; doing so would place the corpus in that system's training data
and invalidates results. It should not be treated as a source of factual
knowledge about any real company (all content is fictional), and reported scores
should always state the tier, the system configuration, and the capability-aware
N/A handling, because aggregate scores are not comparable across different
harness settings.

\textbf{Are there tasks for which the dataset should not be used?}
{\sloppy
It is not a conversational or personal-memory benchmark and should not be
compared head-to-head against LoCoMo/LongMemEval-style results, which measure a
different (and now largely saturated) regime.\par}

%------------------------------------------------------------------------------
\subsection{Distribution}
%------------------------------------------------------------------------------

\textbf{How will the dataset be distributed?}
The full Helix corpus and ground truth are distributed in the public project
repository at \url{https://github.com/JackCGardner/OrgMemBench} under
\texttt{datasets/helix/}, which is the single authoritative copy; all tiers live
in-repo.

\textbf{Under what license is the dataset distributed?}
The \textbf{dataset} is released under \textbf{Creative Commons Attribution 4.0
(CC BY 4.0)} (see \texttt{datasets/LICENSE},
\texttt{SPDX-License-Identifier: CC-BY-4.0}). The accompanying \textbf{evaluation
harness and code} are released under the \textbf{MIT License}.

\textbf{Are there third-party IP or usage restrictions?}
No. The data is entirely synthetic and original to this project; there are no
third-party copyright, license, or usage restrictions on the artifact text, and
no proprietary content is embedded.

\textbf{When will the dataset be distributed?}
The benchmark, harness, and a planned public leaderboard are released with this
paper, on GitHub, to support reproducible progress.

%------------------------------------------------------------------------------
\subsection{Maintenance}
%------------------------------------------------------------------------------

\textbf{Who will maintain the dataset?}
The dataset is maintained by the authors via the project repository. Issues,
errata reports, and questions can be filed through the repository's issue
tracker.

\textbf{Will the dataset be updated, and how will updates be communicated?}
{\sloppy
Yes. The dataset is versioned (\texttt{benchmark\_v0.0.jsonl}), and every
correction is recorded in per-tier \texttt{CHANGELOG.md} files with pre-patch
snapshots preserved. Planned updates include the large-tier and above EMERGENT
questions (the only tiers in which that category is in scope), the XL/XXL scale
tiers, additional companies of other sizes
and verticals for generalization claims, and a cross-model validation pass.
Updates will be communicated through the repository changelog and tagged GitHub
releases.\par}

\textbf{Will older versions continue to be supported?}
Yes. Pre-patch snapshots and version-suffixed benchmark files are retained in the
repository so prior versions remain reproducible, and the changelog documents
what changed between versions.

\textbf{If others want to extend or contribute to the dataset, is there a
mechanism to do so?}
Yes. Because the generation pipeline (\texttt{helix-corpus}, Stages A--F) and the
source-of-truth graph format are part of the release, contributors can generate
additional companies or tiers following the documented stages and submit them
through the repository.
